% from ega4-8-fr.tex
\oldpage[IV-3]{31}
de $f(X)$, et par suite il y a un ouvert quasi-compact $Y'$ dans $Y$ qui
contient $f(X)$~; il existe par suite un $\lambda\geq\alpha$ et un ouvert
quasi-compact $Y'_\lambda$ dans $Y_\lambda$ tels que
$Y'=w_\lambda^{-1}(Y'_\lambda)$ (8.2.11). Remplaçant au besoin $\alpha$ par
$\lambda$ et $Y_\lambda$ par $Y'_\lambda$, on peut donc se borner au cas où
$Y_\alpha$ est quasi-compact, donc aussi $Y$ et les $Y_\lambda$. Comme $Y$ est
quasi-compact, il est réunion finie d'ouverts affines $V_i$, et par suite $X$
est réunion des ouverts $f^{-1}(V_i)$. Comme tout point de $X$ a un voisinage
ouvert quasi-compact contenu dans un des $f^{-1}(V_i)$ et que $X$ est
quasi-compact, on peut, en tenant compte de (8.2.11) et remplaçant au besoin
$\alpha$ par un indice $\lambda\geq\alpha$, supposer que $Y$ est réunion finie
d'ouverts $V_i=w_\alpha^{-1}(V_{i\alpha})$, où les $V_{i\alpha}$ sont des
ouverts affines de $Y_\alpha$~; par suite $X$ est réunion des ouverts
$f^{-1}(V_i)$. Comme tout point de $X$ a un voisinage ouvert quasi-compact
contenu dans l'un des $f^{-1}(V_i)$ et que $X$ est quasi-compact, on peut
recouvrir $X$ par un nombre fini de tels voisinages, et (en répétant au besoin
certains des $V_i$) supposer que l'on a un recouvrement $(U_i)$ de $X$ par des
ouverts quasi-compacts ayant même ensemble d'indices que $(V_i)$ et tel que
$f(U_i)\subset V_i$ pour tout $i$. En outre, à l'aide de (8.2.11) (et en
remplaçant au besoin $\alpha$ par un indice $\lambda\geq\alpha$), on peut
supposer que l'on a $U_i=v_\alpha^{-1}(U_{i\alpha})$ où les $U_{i\alpha}$ sont
des ouverts quasi-compacts dans $X_\alpha$~; de plus, utilisant (8.3.4) comme
ci-dessus, on peut supposer que $(U_{i\alpha})$ est un \emph{recouvrement} de
$X_\alpha$.

Cela étant, il suffira de montrer qu'il existe un $\lambda\geq\alpha$ et pour
chaque $i$ un morphisme $f_{i\lambda}:U_{i\lambda}\to V_{i\lambda}$ (avec
$U_{i\lambda}=v_{\alpha\lambda}^{-1}(U_{i\alpha})$ et
$V_{i\lambda}=w_{\alpha\lambda}^{-1}(V_{i\alpha})$) tels que le morphisme
correspondant $f_i=e_\lambda(f_{i\lambda}):U_i\to V_i$ soit égal à la
\emph{restriction} de $f$ à $U_i$. En effet, s'il en est ainsi, comme
$X_\lambda$ est quasi-séparé (I, 1.2.3), les
$U_{i\lambda}\cap U_{j\lambda}$ sont \emph{quasi-compacts} et le résultat
d'unicité démontré plus haut (qui s'applique puisque $Y_\lambda$ est
localement de type fini sur $S_\lambda$) prouve qu'il existe
$\mu\geq\lambda$ tel que $f_{i\mu}$ et $f_{j\mu}$ coïncident dans
$U_{i\mu}\cap U_{j\mu}$ pour tout couple $(i,j)$, et par suite définissent un
$S_\mu$-morphisme $f_\mu:X_\mu\to Y_\mu$ tel que $e_\mu(f_\mu)=f$.

On est ainsi ramené au cas où $Y_\alpha$ est \emph{affine}, et comme en outre
on peut supposer que les $V_{i\alpha}$ ont des images dans $S_\alpha$
contenues dans des ouverts affines, on peut aussi supposer que $S_\alpha$ est
affine, soient donc $S_\alpha=\operatorname{Spec}(A_\alpha)$,
$Y_\alpha=\operatorname{Spec}(C_\alpha)$, $C_\alpha$ étant une
$A_\alpha$-algèbre de présentation finie, $S=\operatorname{Spec}(A)$,
$Y=\operatorname{Spec}(C)$, avec $A=\varinjlim A_\lambda$,
$C=\varinjlim(C_\alpha\otimes_{A_\alpha}A_\lambda)=
C_\alpha\otimes_{A_\alpha}A$. On a alors
\[
\operatorname{Hom}_S(X,Y)=
\operatorname{Hom}_{A\text{-alg.}}(C,\Gamma(X,\mathcal O_X))=
\operatorname{Hom}_{A_\alpha\text{-alg.}}(C_\alpha,\Gamma(X,\mathcal O_X))
\]
(I, 2.2.4) et de même
\[
\operatorname{Hom}_{S_\lambda}(X_\lambda,Y_\lambda)=
\operatorname{Hom}_{A_\lambda\text{-alg.}}
(C_\alpha\otimes_{A_\alpha}A_\lambda,\Gamma(X_\lambda,\mathcal O_{X_\lambda}))=
\operatorname{Hom}_{A_\alpha\text{-alg.}}
(C_\alpha,\Gamma(X_\lambda,\mathcal O_{X_\lambda})).
\]
Mais comme $X_\alpha$ est quasi-compact et quasi-séparé, on sait (8.5.4) que
l'on a
$\varinjlim\Gamma(X_\lambda,\mathcal O_{X_\lambda})=
\Gamma(X,\mathcal O_X)$~; puisque $C_\alpha$ est une $A_\alpha$-algèbre de
présentation finie, le fait que (8.8.1.1) est bijectif résulte alors de
(8.8.2.3).

Avant de passer à la démonstration de (ii), notons les corollaires suivants de
(i)~:
\end{proof}

\begin{corollary}[8.8.2.4]
\label{IV.8.8.2.4-fr}
Supposons $S_0$ quasi-compact, $X_\alpha$ de présentation finie sur $S_\alpha$,
$Y_\alpha$ de type fini sur $S_\alpha$ et quasi-séparé sur $S_\alpha$ (ce qui
sera par exemple le cas si $Y_\alpha$ est également
